% !TeX program = lualatex
% !TeX encoding = utf8
% !TeX spellcheck = uk_UA
% !TeX root = ../main.tex

%% --------------------------------------------------------
%% Chapter file: Scales.tex
%% --------------------------------------------------------

\graphicspath{{\currfilebase/Pictures}}

\chapter{Основні масштаби, режими та межі моделей}
\label{sec:plasma-scales}\hypertarget{plasma-scales}{}
%% --------------------------------------------------------

Після розгляду основних фізичних явищ корисно зібрати в одному місці
методичну ідею, яка проходить через усю фізику плазми: \emphx{модель
визначається не назвою середовища, а співвідношенням характерних
масштабів}. Одна й та сама плазма може потребувати різного опису
залежно від просторового масштабу, характерного часу та швидкості
процесу.

Нехай $L$ --- просторовий масштаб задачі, а
$\tau\sim1/\omega$ --- характерний час її розвитку. Насамперед
порівнюють

\begin{equation*}
 \frac{\lambda_D}{L},\qquad
 \frac{r_L}{L},\qquad
 \frac{\omega}{\omega_p},\qquad
 \frac{\omega}{\omega_c},\qquad
 \frac{\nu}{\omega_c}.
\end{equation*}

Ці відношення не є незалежним набором \enquote{характеристик плазми}. Кожне
з них відповідає на конкретне фізичне питання: чи можна вважати плазму
квазінейтральною, чи замагнічені частинки, чи встигають електрони
відновлювати квазінейтральність, чи встигають частинки здійснювати
гірообертання та чи руйнують зіткнення замагніченість.

%% ========================================================
\section{Просторові та часові масштаби}
%% ========================================================

Для токамака ($n=10^{14}~\mathrm{см}^{-3}$, $T=10$~кеВ, $B=5\times10^{4}$~Гс, $L\approx1$~м)
ієрархію довжин і~частот показано на рис.~\ref{fig:scales-hierarchy}; далі кожна умова вибору
моделі~--- це порівняння двох точок на цій шкалі.

\begin{figure}[ht]\centering
	\localinput{scales-hierarchy.tikz}
	\caption{Ієрархія масштабів у~токамаку на логарифмічній осі: довжини (зверху) і~частоти (знизу).
	Червоним позначено масштаби задачі ($L$, $\lambda_{\mathrm{пр}}$, $\nu_{ei}$), синім~--- магнітні масштаби.
	Для $\omega_{LH}$ використано точну формулу з~урахуванням $\omega_{pi}$}
	\label{fig:scales-hierarchy}
\end{figure}

Умова

\begin{equation*}
\lambda_D\ll L
\end{equation*}

означає, що дебаївське екранування відбувається на значно меншому
масштабі, ніж досліджувана задача. Тому на масштабі $L$ можна, як
правило, покладати $n_e\simeq Zn_i$ і користуватися
квазінейтральним описом. Якщо ж $L\sim\lambda_D$, електростатичне поле
та розділення зарядів стають частиною самої задачі.

Аналогічно,

\begin{equation*}
r_L\ll L
\end{equation*}

є просторовою умовою: на одній ларморівській орбіті поле та його градієнти
майже однорідні, тому орбіту можна замінити ведучим центром. Вона відрізняється від
часової умови $\omega\ll\omega_c$ (за час зміни плазми частинка багато разів
обертається навколо силової лінії); обидві потрібні для дрейфового наближення
(розд.~\ref{sec:particle-motion}). Якщо $r_L\sim L$, гірація вже
не є малим швидким масштабом, який можна відокремити від основного
руху.

Порівняння характерної частоти процесу з плазмовою частотою має інший
зміст. Умова $\omega\ll\omega_p$ означає, що електрони швидко реагують
на спробу порушити квазінейтральність, тоді як при
$\omega\gtrsim\omega_p$ їхня інерція стає суттєвою.

Умова $\omega\ll\omega_c$ означає, що гірація є швидким рухом
порівняно з еволюцією макроскопічного збурення. Саме тому у повільній
МГД-динаміці окремий гірооберт можна не відстежувати.

Нарешті, відношення $\nu/\omega_c$ характеризує конкуренцію зіткнень
та магнітного обертання. При $\nu\ll\omega_c$ частинка встигає
здійснити багато гірообертів між зіткненнями і є сильно замагніченою;
при $\nu\gg\omega_c$ зіткнення руйнують гірообертання ще до завершення
циклу.

%% ========================================================
\section{Зіткнювальний та беззіткнювальний режими}
%% ========================================================

Для зіткнювального та беззіткнювального режимів одного порівняння
частот недостатньо. Вводимо середню довжину вільного пробігу

\begin{equation*}
\lambda_{\mathrm{пр}}\sim\frac{v_{\mathrm{th}}}{\nu}
\end{equation*}

та число Кнудсена

\begin{equation}
Kn=\frac{\lambda_{\mathrm{пр}}}{L}.
\label{eq:plasma-knudsen}
\end{equation}

Якщо $Kn\ll1$, частинка багато разів стикається на масштабі $L$ і
локальний гідродинамічний опис може бути виправданим. Якщо
$Kn\gtrsim1$, структура функції розподілу за швидкостями стає
істотною, і кінетичні ефекти вже не можна автоматично відкидати.

\emph{Ключова оцінка для токамака} ($n=10^{14}~\mathrm{см}^{-3}$, $T=10$~кеВ): $\nu_{ei}\approx5\times10^{3}~\mathrm{с}^{-1}$, $v_{Te}\approx4\times10^{9}~\mathrm{см/с}$,
тому $\lambda_{\mathrm{пр}}=v_{Te}/\nu_{ei}\approx10^6$~см $\approx10$~км, тоді як $L\approx1$~м. Отже, $Kn\sim10^4$: плазма \enquote{беззіткнювальна},
проте МГД працює. Парадокс знімається так: поперек поля роль довжини вільного пробігу відіграє $r_{Li}\ll L$ (магнітне поле \enquote{склеює} частинки поперек $\vect B$),
а~вздовж поля ні~--- тому саме поздовжній перенос потребує кінетики.

Водночас беззіткнювальність завжди треба розуміти відносно
конкретного процесу. Умова

\begin{equation*}
\nu\ll\omega
\end{equation*}

означає, що зіткнення слабкі саме для явища з характерною частотою
$\omega$. Повільний процес може бути зіткнювальним у тій самій плазмі,
для якої швидкі коливання є беззіткнювальними.

Отже, твердження \enquote{плазма беззіткнювальна} без зазначення масштабу
часу є неповним.

%% ========================================================
\section{Як за масштабами вибирають модель}
%% ========================================================

Корисно розглядати наведені вище співвідношення не окремо, а як
послідовність запитань.

\begin{table}[ht]
	\centering\small
	\caption{Характерні умови вибору моделі плазми.}
	\label{tab:plasma-regimes}
	\begin{platbl}{colspec={X[l,m]X[c,m]X[l,m]}}
Умова & Фізичний наслідок & Типовий опис \\
$\lambda_D\ll L$ & квазінейтральність, $n_e\simeq Zn_i$ & можна вводити єдину рідину (потрібне ще замикання) \\
$r_L\ll L$ & замагніченість & дрейфовий / МГД (разом з~$\omega\ll\omega_{ci}$) \\
$Kn\ll1$ & багато зіткнень на масштабі $L$ & гідродинамічний \\
$Kn\gtrsim1$ & важлива структура $f(\vect r,\vect v,t)$ & кінетичний \\
$\nu\ll\omega$ & зіткнення слабкі для даного процесу & беззіткнювальний режим \\
$R_m\gg1$ & резистивна дифузія поля повільна & ідеалізована МГД \\
$\beta\ll1$ & магнітний тиск домінує & магнітно-керований режим \\
$\beta\gtrsim1$ & плазма істотно деформує поле & МГД з~повним балансом тиску плазми й~поля \\
	\end{platbl}
\end{table}

Таблицю~\ref{tab:plasma-regimes} не слід читати як набір
взаємовиключних категорій. Реальна плазма може одночасно задовольняти
кілька нерівностей. Саме їхня сукупність визначає, які члени рівнянь
можна відкинути, а які залишити.

Це особливо важливо для МГД. Наприклад, $R_m\gg1$ саме по собі не
означає повної ідеальної МГД: воно лише пригнічує резистивний член.
Щоб отримати

\begin{equation*}
\vect E+\frac{1}{c}\vect v\times\vect B\simeq0,
\end{equation*}

треба також перевірити малість холлівського, електроннотискового та
інерційного членів узагальненого закону Ома.

Таким чином, \enquote{МГД-опис} --- це не властивість плазми взагалі, а~наближення, справедливість якого треба обґрунтувати для конкретного
процесу.

%% ========================================================
\section{Типова ієрархія масштабів у магнітному утриманні}
%% ========================================================

Розглянемо типову установку магнітного утримання:
$n=10^{14}~\mathrm{см}^{-3}$, $T=10$~кеВ,
$B=5\times10^{4}$~Гс, дейтерієва плазма. Для неї характерна ієрархія

\begin{equation*}
\underbrace{n^{-1/3}}_{2\times10^{-5}~\text{см}}
\ll
\underbrace{\lambda_{De}}_{7\times10^{-3}~\text{см}}
\sim
\underbrace{r_{Le}}_{5\times10^{-3}~\text{см}}
\ll
\underbrace{r_{Li}}_{0{,}3~\text{см}}
\ll
\underbrace{L}_{\sim10^{2}~\text{см}}.
\end{equation*}

Ця послідовність дає набагато більше фізичної інформації, ніж кожне
число окремо. На найменшому масштабі видно дискретність частинок;
далі виникають дебаївське екранування та електронна гірація; ще більшим
є іонний ларморівський радіус; і лише на масштабі всієї установки
з'являється макроскопічний рух плазми.

Частотна ієрархія для тієї самої оцінки має вигляд

\begin{equation*}
\begin{gathered}
\omega_{pe}\approx6\times10^{11}~\mathrm{с}^{-1},
\qquad
\omega_{ce}\approx9\times10^{11}~\mathrm{с}^{-1},\\
\omega_{ci}\approx2\times10^{8}~\mathrm{с}^{-1},
\qquad
\nu_{ei}\approx5\times10^{3}~\mathrm{с}^{-1}.
\end{gathered}
\end{equation*}

Отже, зіткнення значно рідші за електронний та іонний циклотронний
рух. Плазма сильно замагнічена і для багатьох швидких процесів
практично беззіткнювальна.

Але з цієї ж ієрархії випливає важливе обмеження: \enquote{беззіткнювальна}
не означає \enquote{завжди кінетична}. На достатньо великих масштабах і~повільних часах багато мікроскопічних ступенів свободи усереднюються,
і виникає ефективний рідинний опис.

%% ========================================================
\section{Від кінетики до МГД}
%% ========================================================

На масштабах, що значно перевищують іонний ларморівський радіус,
і для повільних рухів, коли $\omega\ll\omega_{ci}$, доцільний
магнітогідродинамічний опис: плазма розглядається як провідне
суцільне середовище (вивід~--- розд.~\ref{sec:kinetic-moments} і~\ref{sec:mhd}).

Для швидких процесів і малих масштабів потрібен кінетичний опис:
функція розподілу

\begin{equation*}
f_s=f_s(\vect r,\vect v,t)
\end{equation*}

залежить від координат і швидкостей. Для беззіткнювальної плазми вона
задовольняє рівняння Власова

\begin{equation}
\frac{\partial f_s}{\partial t}
+\vect v\cdot\frac{\partial f_s}{\partial\vect r}
+\frac{q_s}{m_s}
\left(
\vect E+\frac{1}{c}\vect v\times\vect B
\right)
\cdot\frac{\partial f_s}{\partial\vect v}
=0.
\label{eq:plasma-vlasov-summary}
\end{equation}

Рівняння~\eqref{eq:plasma-vlasov-summary} доповнюється рівняннями
Максвелла з~джерелами

\begin{equation*}
\rho=\sum_s q_s\int f_s\,d^3v,
\qquad
\vect j=\sum_s q_s\int\vect v f_s\,d^3v.
\end{equation*}

Саме з кінетичного опису природно випливають ефекти, які не містяться
в~найпростішій МГД: загасання Ландау, кінетичні плазмові хвилі,
резонансна взаємодія хвиль із частинками та багато видів
нестійкостей.

Водночас кінетичний опис не \enquote{кращий} за МГД у всіх ситуаціях.
Він містить значно більше інформації, ніж потрібно для повільної
макроскопічної динаміки. Завдання теорії --- не використовувати
найскладнішу модель, а відкинути лише ті ступені свободи, які
справді є несуттєвими на вибраному масштабі.

%% ========================================================
\section{Безрозмірні параметри МГД}
%% ========================================================

Для магнітогідродинаміки додатково важливі безрозмірні параметри

\begin{equation}
\beta=\frac{8\pi p}{B^2},
\qquad
R_m\sim\frac{4\pi\sigma Lv}{c^2},
\qquad
M_A=\frac{v}{v_A},
\qquad
v_A=\frac{B}{\sqrt{4\pi\rho_m}}.
\label{eq:plasma-mhd-parameters}
\end{equation}

Параметр $\beta$ порівнює плазмовий тиск із магнітним. При
$\beta\ll1$ магнітне поле домінує над газодинамічним тиском, тоді як
при $\beta\gtrsim1$ тиск плазми стає істотним у силовому балансі. Сама величина $\beta$ не~визначає
застосовності МГД (вона працює і~за малих $\beta$), а~лише показує, наскільки плазма
впливає на~поле.

Магнітне число Рейнольдса $R_m$ порівнює характерний час руху плазми
з часом дифузії магнітного поля. Велике $R_m$ означає, що резистивна
дифузія поля повільна порівняно з рухом плазми.

Число Альвена $M_A$ порівнює швидкість руху плазми з~альвенівською
швидкістю. Воно показує, наскільки важливі магнітні напруження для
динаміки порівняно з~інерцією потоку.

Ці параметри не задають одну-єдину \enquote{категорію плазми}. Вони описують
різні сторони одного й того самого режиму.

%% ========================================================
\section{Межі простих моделей}
%% ========================================================

Наведені вище критерії дозволяють визначити режим, але не скасовують
обмежень самих моделей.

У цьому курсі ми лише позначили, але не розвинули повністю
дрейфово-кінетичну теорію, нелінійну теорію хвиль та сильну
турбулентність. Повний опис хвиль у~замагніченій плазмі потребує розв'язання
тензорного дисперсійного рівняння, а~детальний опис нестійкостей ---
лінеаризації кінетичних або МГД-рівнянь для конкретної рівноваги.

Особливо важливо, що реальний перенос у магнітно утримуваній плазмі
може бути значно сильнішим за класичний колізійний перенос.
Турбулентність породжує \emph{аномальний перенос}, який може визначати
реальний час утримання енергії. Тому навіть дуже мала частота
зіткнень не гарантує ідеального утримання. Оцінка: для токамака бомівська дифузія $D_B=cT/16eB\approx1{,}3\times10^{6}~\mathrm{см}^2/\mathrm{с}$ ($\approx10^{2}~\mathrm{м}^2/\mathrm{с}$),
тоді як класична $D_\perp\sim\nu_{ei}r_{Le}^2\sim0{,}1$~см$^2$/с: розрив у~$\sim10^{7}$ разів. Турбулентні флуктуації $\vect E\times\vect B$ перемішують плазму
з~кроком порядку розміру вихору, а~не $r_L$ (розд.~\ref{sec:magnetic-confinement}).

Сильнозв'язана та квантова плазма потребує окремих методів. Отже,
опис, побудований у цій книзі, треба розуміти як систему моделей із
чітко визначеними межами застосовності, а не як універсальне рівняння
для будь-якої плазми.

%% ========================================================
\section{Основні масштаби та формули}
%% ========================================================

У таблиці зведено основні величини посібника. Наведені практичні
формули відповідають системі СГС; $T$ задано в~еВ, $n$ --- у~$\mathrm{см}^{-3}$, $B$ --- у Гс, $f$ --- у МГц; $A$ і~$Z$ ---
масове число та заряд іона.

\begin{Attention}
	Практичні формули в~таблиці записано в~СГС із~числовими множниками для вказаних одиниць ($T$ в~еВ, $n$ в~см$^{-3}$, $B$ в~Гс); вони не~переносяться на~СІ без перерахунку.
	Температуру всюди виражено в~енергетичних одиницях (еВ).
\end{Attention}

{\small
\begin{longtblr}[caption={Основні масштаби та формули фізики плазми в системі СГС.},label={tab:plasma-formulas}]{colspec={X[l,m]X[c,m]X[l,m]}}
		Величина & Формула & Практична форма \\
		Дебаївська довжина &
		$\lambda_{De}=\sqrt{T_e/4\pi n_ee^2}$ &
		$7{,}43\times10^{2}\sqrt{T_e/n_e}$~см \\
		Число частинок у сфері &
		$N_D=\frac{4\pi}3n_e\lambda_{De}^3$ &
		$N_D\approx0{,}094\,\Gamma^{-3/2}$ \\
		Параметр зв'язку &
		$\Gamma=e^2n^{1/3}/T$ & \\
		Плазмова частота &
		$\omega_{pe}=\sqrt{4\pi n_ee^2/m_e}$ &
		$5{,}64\times10^{4}\sqrt{n_e}~\mathrm{с}^{-1}$ \\
		Теплова швидкість &
		$v_{Te}=\sqrt{T_e/m_e}$ &
		$4{,}19\times10^{7}\sqrt{T_e}$~см/с \\
		Циклотронна частота &
		$\omega_c=|q|B/mc$ &
		$\omega_{ce}=1{,}76\times10^{7}B$,
		$\omega_{ci}=9{,}58\times10^{3}(Z/A)B$ \\
		Ларморівський радіус &
		$r_L=mcv_\perp/|q|B$ &
		$r_{Le}=2{,}38\sqrt{T_e}/B$~см,
		$r_{Li}=1{,}02\times10^{2}\sqrt{AT_i}/(ZB)$~см \\
		Частота зіткнень $e$--$i$ &
		формула~\eqref{eq:nu-ei} &
		$2{,}9\times10^{-6}\,n\ln\Lambda\,T_e^{-3/2}~\mathrm{с}^{-1}$ \\
		Провідність (Спітцер), вздовж $\vect B$ &
		$\sigma_\parallel\propto T_e^{3/2}/\ln\Lambda$ &
		$1{,}7\times10^{14}\,T_e^{3/2}/\ln\Lambda~\mathrm{с}^{-1}$ (поперек $\sigma_\perp\approx\sigma_\parallel/2$) \\
		Критична концентрація &
		$n_{\mathrm{кр}}=m_e\omega^2/4\pi e^2$ &
		$1{,}24\times10^{4}f^2~\mathrm{см}^{-3}$ \\
		Магнітний тиск &
		$B^2/8\pi$ &
		$3{,}98\times10^{6}\,(B/10^4)^2$~дин/$\mathrm{см}^2$ \\
		Параметр $\beta$ &
		$\beta=8\pi p/B^2$ & \\
		Альвенівська швидкість &
		$v_A=B/\sqrt{4\pi\rho_m}$ &
		$2{,}18\times10^{11}B/\sqrt{n_i}$~см/с (водень) \\
		Час дифузії поля &
		$\tau_d=4\pi\sigma L^2/c^2$ & \\
		Довжина Томаса --- Фермі &
		$\lambda_{TF}=\sqrt{E_F/6\pi ne^2}$ &
		$0{,}5\sqrt{a_Bn^{-1/3}}$ \\
		Амбіполярна дифузія &
		$D_a\approx D_i(1+T_e/T_i)$ & \\
		Швидкість іонного звуку &
		$c_s=\sqrt{(ZT_e+3T_i)/m_i}$ &
		$9{,}79\times10^{5}\sqrt{ZT_e/A}$~см/с ($T_i\ll T_e$) \\
		Дрейф $\vect E\times\vect B$ &
		$v_E=cE/B$ &
		$10^{8}\,E[\text{В/см}]/B[\text{Гс}]$~см/с \\
		Дифузія Бома &
		$D_B=cT_e/16eB$ &
		$6{,}25\times10^{6}\,T_e/B~\mathrm{см}^2/\mathrm{с}$ \\
		Верхній / нижній гібридні частоти &
		$\omega_{UH}^2=\omega_{pe}^2+\omega_{ce}^2$; $\omega_{LH}^2\simeq\omega_{ci}\omega_{ce}$ (лише при $\omega_{pe}\gg\omega_{ce}$) &
		$f_{UH}=\sqrt{f_{pe}^2+f_{ce}^2}$ \\
		Поле Драйсера &
		$E_D=4\pi e^3n\ln\Lambda/T_e$ &
		$2{,}6\times10^{-13}\,n\ln\Lambda/T_e~\text{В/см}$ \\
		Критерій запалювання (D--T) &
		$n\tau_E\geq12T/\langle\sigma v\rangle E_\alpha$ &
		$\gtrsim1{,}5\times10^{14}~\mathrm{см}^{-3}\mathrm{с}$ (мінімум, $T\sim25$~кеВ) \\
		Гальмівне випромінювання &
		$P_{br}\propto n_e\sqrt{T_e}\sum Z^2n_i$ &
		$1{,}7\times10^{-32}n_e\sqrt{T_e}\sum Z^2n_i$ Вт/$\mathrm{см}^3$ \\
	\end{longtblr}}

Таблиця~\ref{tab:plasma-formulas} призначена не для заміни виведень,
а для швидкої оцінки порядків величин. При виборі моделі спочатку
визначають характерні масштаби задачі, а вже потім звертаються до
відповідних формул.

%% ========================================================
\section{Що вимірюють на~якому масштабі}\label{sec:scales-diagnostics}
%% ========================================================

Кожен параметр плазми \enquote{видно} лише методом, чий власний масштаб (довжина хвилі, ларморівський радіус, розмір котушки) сумірний із~вимірюваною величиною.
Таблиця~\ref{tab:scales-methods} зводить методи, розібрані в~попередніх розділах.

\begin{table}[ht]
	\caption{Масштаб, величина й~метод діагностики.}
	\label{tab:scales-methods}
	\centering\small
	\begin{platbl}{colspec={X[l,m]X[l,m]X[l,m]}}
Величина & Метод & Де розглянуто \\
$n_e$ (інтегральна, по~хорді) & інтерферометрія, $\omega\gg\omega_p$ & розд.~\ref{sec:wave-diagnostics} \\
$n_e$ (локальна, профіль) & рефлектометрія, $\omega=\omega_p$ & розд.~\ref{sec:wave-diagnostics} \\
$T_e$, $n_e$ (локально) & томсонівське розсіяння, $\alpha=1/k\lambda_D$ & розд.~\ref{sec:thomson} \\
$T_e$ (середня) & опір за~Спітцером & розд.~\ref{sec:collision-diagnostics} \\
$T_e(r)$ & ЕЦЕ, $\omega=\ell\omega_{ce}$ & розд.~\ref{sec:motion-diagnostics} \\
$T_i$ & допплерівське розширення ліній, нейтрони & розд.~\ref{sec:collision-diagnostics}, \ref{sec:fusion-diagnostics} \\
$I$, енергія плазми & пояс Роговського, діамагнітна петля & розд.~\ref{sec:magnetic-diagnostics} \\
положення, $\Delta$, $\beta_p+l_i/2$ & магнітна реконструкція & розд.~\ref{sec:equilibrium-diagnostics} \\
моди $m,n$, $\delta B$ & котушки Мірнова, м'який рентген & розд.~\ref{sec:instability-diagnostics} \\
$q(r)$, $B$ & обертання Фарадея & розд.~\ref{sec:faraday-diagnostics} \\
$P_f$, $Q$ & потік нейтронів & розд.~\ref{sec:fusion-diagnostics} \\
	\end{platbl}
\end{table}

Метод треба вибирати за~тими самими критеріями, що й~модель: наприклад, томсонівське розсіяння з~$\alpha\ll1$ дає тепловий спектр, а~з~$\alpha\gtrsim1$~--- колективні ефекти.

\begin{Summary}
\begin{itemize}
\item Модель плазми визначається не назвою середовища, а~співвідношенням характерних масштабів задачі.
\item Порівнюють $\lambda_D/L$, $r_L/L$, число Кнудсена $Kn$, частоти $\nu/\omega$, магнітне число Рейнольдса $R_m$ та~$\beta$.
\item Зіткнювальний режим допускає гідродинамічний опис, беззіткнювальний~--- потребує кінетики.
\item Простіші моделі мають межі застосовності; їх треба перевіряти оцінкою масштабів до розв'язання задачі.
\item Параметри плазми вимірюють методами, чий масштаб відповідає вимірюваній величині (табл.~\ref{tab:scales-methods}).
\end{itemize}
\end{Summary}

\Needspace{9\baselineskip}
\section*{Задачі}
%% --------------------------------------------------------

\begin{problem}\label{prb:PlasmaTable}
Оцініть $\lambda_{De}$, середню відстань
$a=n^{-1/3}$ та $N_D$ для чотирьох об'єктів з~табл.~\ref{tab:plasma-types}:
іоносфера ($n=10^5~\mathrm{см}^{-3}$, $0{,}1$~еВ), сонячна корона
($10^8~\mathrm{см}^{-3}$, $100$~еВ), токамак
($10^{14}~\mathrm{см}^{-3}$, $10$~кеВ), ядро Сонця
($10^{26}~\mathrm{см}^{-3}$, $1$~кеВ). Для якого з них колективний
опис найменш обґрунтований?
\emph{Відповідь:} іоносфера: $\lambda_D\approx0{,}7$~см, $a\approx2\times10^{-2}$~см, $N_D\approx2\times10^{5}$;
корона: $0{,}7$~см, $2\times10^{-3}$~см, $2\times10^{8}$; токамак: $7\times10^{-3}$~см, $2\times10^{-5}$~см, $2\times10^{8}$;
ядро Сонця: $\lambda_D\approx2\times10^{-9}$~см${}\sim a$, $N_D\approx5$~--- колективний опис найменш обґрунтований.
\end{problem}

\begin{problem}\label{prb:PlasmaScales}
Для дейтерієвої плазми $n=10^{14}~\mathrm{см}^{-3}$,
$T=10$~кеВ у~полі $B=5\times10^{4}$~Гс знайдіть
$n^{-1/3}$, $\lambda_{De}$, $r_{Le}$, $r_{Li}$ і розташуйте їх за
зростанням. Яка з обох довжин ($r_{Le}$ чи $\lambda_{De}$) більша?
\emph{Відповідь:} $n^{-1/3}\approx2\times10^{-5}$~см${}<r_{Le}\approx4{,}8\times10^{-3}$~см${}<\lambda_{De}\approx7{,}4\times10^{-3}$~см${}<r_{Li}\approx0{,}29$~см;
$\lambda_{De}$ трохи більша за~$r_{Le}$.
\end{problem}

\begin{problem}\label{prb:PlasmaCollisions}
Для $n_e=10^{12}~\mathrm{см}^{-3}$, $T_e=10$~еВ,
$\ln\Lambda=15$ знайдіть частоту зіткнень $\nu_{ei}$, відношення
$\nu_{ei}/\omega_{pe}$ та довжину вільного пробігу
$\lambda_{\mathrm{пр}}=v_{Te}/\nu_{ei}$. Порівняйте
$\lambda_{\mathrm{пр}}$ з $\lambda_{De}$.
\emph{Відповідь:} $\nu_{ei}\approx1{,}4\times10^{6}~\mathrm{с}^{-1}$, $\nu_{ei}/\omega_{pe}\approx2{,}5\times10^{-5}$,
$\lambda_{\mathrm{пр}}\approx1$~м; $\lambda_{\mathrm{пр}}/\lambda_{De}\approx4\times10^{4}$.
\end{problem}

\begin{problem}\label{prb:PlasmaFormulaCheck}
Користуючись табл.~\ref{tab:plasma-formulas}, оцініть $\lambda_D$
та $\omega_{pe}$ для плазми газового розряду з~$n_e=10^{10}~\mathrm{см}^{-3}$ і~$T_e=3$~еВ. Порівняйте отримані масштаби з типовим
часом і просторовим масштабом макроскопічного процесу.
\emph{Відповідь:} $\lambda_D\approx1{,}3\times10^{-2}$~см, $\omega_{pe}\approx5{,}6\times10^{9}~\mathrm{с}^{-1}$ ($f\approx0{,}9$~ГГц);
обидва масштаби набагато менші за~сантиметрові розміри й~мікро-/мілісекундні часи розряду, тому плазма квазінейтральна.
\end{problem}

\begin{problem}\label{prb:ClassicalVsBohmTau}
Оцініть час утримання $\tau\sim a^2/D$ для токамака з~$a=1$~м і~$B=5$~Тл, $T=10$~кеВ, якщо
діє (а)~класична дифузія $D_\perp\sim0{,}1~\mathrm{см}^2/\mathrm{с}$, (б)~бомівська $D_B=cT/(16eB)$.
Який із~двох результатів узгоджується з~експериментом ($\tau_E\sim1$~с) і~що це означає для вибору моделі?
\emph{Відповідь:} (а)~$\sim10^{5}$~с, (б)~$\sim10^{-2}$~с.
\end{problem}

\begin{problem}\label{prb:IdealDegenerate}
За якої густини плазма з~$T=1$~кеВ перестає бути ідеальною ($\Gamma\sim1$) і~за якої стає виродженою
($T\sim E_F$)? Порівняйте з~ядром Сонця (задача~\ref{prb:PlasmaTable}).
\emph{Відповідь:} $\Gamma\sim1$ при $n\sim3\times10^{29}~\mathrm{см}^{-3}$, $T\sim E_F$ при $n\sim10^{26}~\mathrm{см}^{-3}$.
\end{problem}

\begin{problem}\label{prb:PlasmaModelLimits}
Для плазми з~$n=10^{14}~\mathrm{см}^{-3}$,
$T_i=10$~кеВ та характерним розміром $L=1$~м задано
$\lambda_{\mathrm{пр}}=10^6~\mathrm{см}$ і $r_{Li}=0{,}3~\mathrm{см}$.
Оцініть $Kn=\lambda_{\mathrm{пр}}/L$ та $r_{Li}/L$. Який із~двох
параметрів показує, що звичайна ізотропна МГД-модель має обмеження?
Яке уточнення опису слід розглядати?
\emph{Відповідь:} $Kn\approx10^{4}$, $r_{Li}/L\approx3\times10^{-3}$; обмеження показує $Kn$: поздовжній перенос
беззіткнювальний, тиск анізотропний~--- потрібні модель CGL, дрейфово- чи гірокінетичний опис поздовжнього переносу.
\end{problem}

\Needspace{8\baselineskip}
\begin{Questions}
\begin{quizlist}

\item Що таке квазінейтральність? Якою нерівністю визначається
можливість нехтувати розділенням зарядів на масштабі $L$?

\item Чому порівняння масштабів часто важливіше за числове значення
окремої величини?

\item Що фізично означають умови $\lambda_D\ll L$ та $r_L\ll L$?

\item Чому умова $\nu\ll\omega$ повинна формулюватися відносно
конкретного процесу?

\item Чим принципово відрізняються МГД- та кінетичний описи?
За яких умов кінетичний опис стає необхідним?

\item Чому виконання $R_m\gg1$ саме по собі не гарантує повної
ідеальної МГД?

\item Як пов'язані між собою число Кнудсена $Kn$, довжина вільного
пробігу та характерний розмір задачі?

\item Чому плазма може бути майже беззіткнювальною для швидких
процесів і водночас поводитися як суцільне середовище на великих
просторових і часових масштабах?

\item Які явища не описуються найпростішою ізотропною МГД-моделлю?

\item Яку інформацію дає типова ієрархія
$n^{-1/3}\ll\lambda_D\sim r_{Le}\ll r_{Li}\ll L$?

\item Чому для МГД важлива саме $\omega_{ci}$, а~не $\omega_{ce}$?

\item Чому в~токамаку перенос поперек поля аномальний, а~не класичний, і~що це означає
для вибору моделі?

\end{quizlist}
\end{Questions}
